% GENERATED FILE. Edit canonical lessons, then run npm run generate:books.
% canonical-source-hash: fnv1a64:5a1dbe41b41c39ae
% canonical-lessons: MR-C357-prasasan, MR-C357-samsad, MR-C357-mantri, MR-C357-neta, MR-C357-nemne, MR-R357-review-350-357

\chapter{Who Governs}
\label{ch:mr-who-governs}
\hlchaptermodality{357}

\begin{chapteropening}
\textbf{By the end of this chapter:} \emph{I can talk about parliament, ministers and leaders.}

Complete the last lesson of chapter 357: I can talk about parliament, ministers and leaders.
\end{chapteropening}

\section[\texorpdfstring{\mr{प्रशासन} (praśāsan)}{प्रशासन (praśāsan)}]{\mr{प्रशासन} (praśāsan) — administration}

\label{lesson:MR-C357-prasasan}



\begin{quote}
Before the new one: say the Marathi for public opinion, then the Marathi for to give a mandate.
\end{quote}

\subsection*{You'll want to know: \mr{प्रशासन}}
\textbf{\mr{प्रशासन}} — \emph{praśāsan} — \textquotedblleft{}administration\textquotedblright{}.

\begin{quote}
\mr{पुरानंतर} \mr{प्रशासनाने} \mr{लोकांना} \mr{मदत} \mr{केली}.
\end{quote}

\begin{quote}
\emph{purānantar praśāsnāne lokānnā madat kelī.}
\end{quote}

\begin{quote}
After the flood the administration helped people.
\end{quote}

Neuter.

Neuter; \mr{प्रशासकीय}, administrative.

\begin{cousinweb}
\textbf{\mr{प्रशासन}} (\emph{praśāsan}) is a tatsama, taken from Sanskrit \emph{praśāsana}, \textquotedblleft{}instruction, government\textquotedblright{} (Monier-Williams).
\end{cousinweb}

\begin{grammarlens}[title={What you've built}]
One more word for giving reasons and opinions.
\end{grammarlens}

\subsection*{Your turn}
\begin{itemize}
\raggedright
  \item \emph{Say it:} \emph{praśāsan}
  \item \emph{Say it:} \emph{praśāsan}, once more
  \item \emph{Say it:} \emph{purānantar praśāsnāne lokānnā madat kelī.} — After the flood the administration helped people.
  \item \emph{Recall:} say the Marathi for public opinion, then the Marathi for to give a mandate, then say \emph{praśāsan} again
\end{itemize}

\begin{tcolorbox}[breakable,colback=teal!4,colframe=teal!35!black,title={Before you move on}]
What does \mr{प्रशासन} mean? (\textquotedblleft{}Administration\textquotedblright{}.) Say it once more.
\end{tcolorbox}
\section[\texorpdfstring{\mr{संसद} (saṁsad)}{संसद (saṁsad)}]{\mr{संसद} (saṁsad) — parliament}

\label{lesson:MR-C357-samsad}



\begin{quote}
Before the new one: say the Marathi for to give a mandate, then the Marathi for administration.
\end{quote}

\subsection*{You'll want to know: \mr{संसद}}
\textbf{\mr{संसद}} — \emph{saṁsad} — \textquotedblleft{}parliament\textquotedblright{}.

\begin{quote}
\mr{संसदेत} \mr{नव्या} \mr{कायद्यावर} \mr{चर्चा} \mr{झाली}.
\end{quote}

\begin{quote}
\emph{saṁsdet navyā kāyadyāvar carcā jhālī.}
\end{quote}

\begin{quote}
The new law was discussed in parliament.
\end{quote}

Feminine.

Feminine; its oblique is \mr{संसदे}-: \mr{संसदेत}, in parliament.

\begin{cousinweb}
\textbf{\mr{संसद}} (\emph{saṁsad}) is a tatsama, taken from Sanskrit \emph{saṁsad}, \textquotedblleft{}an assembly, a session\textquotedblright{} (Monier-Williams).
\end{cousinweb}

\begin{grammarlens}[title={What you've built}]
One more word for giving reasons and opinions.
\end{grammarlens}

\subsection*{Your turn}
\begin{itemize}
\raggedright
  \item \emph{Say it:} \emph{saṁsad}
  \item \emph{Say it:} \emph{saṁsad}, once more
  \item \emph{Say it:} \emph{saṁsdet navyā kāyadyāvar carcā jhālī.} — The new law was discussed in parliament.
  \item \emph{Recall:} say the Marathi for to give a mandate, then the Marathi for administration, then say \emph{saṁsad} again
\end{itemize}

\begin{tcolorbox}[breakable,colback=teal!4,colframe=teal!35!black,title={Before you move on}]
What does \mr{संसद} mean? (\textquotedblleft{}Parliament\textquotedblright{}.) Say it once more.
\end{tcolorbox}
\section[\texorpdfstring{\mr{मंत्री} (mantrī)}{मंत्री (mantrī)}]{\mr{मंत्री} (mantrī) — a cabinet minister; a minister}

\label{lesson:MR-C357-mantri}



\begin{quote}
Before the new one: say the Marathi for administration, then the Marathi for parliament.
\end{quote}

\subsection*{You'll want to know: \mr{मंत्री}}
\textbf{\mr{मंत्री}} — \emph{mantrī} — \textquotedblleft{}a cabinet minister; a minister\textquotedblright{}.

\begin{quote}
\mr{मंत्र्यांनी} \mr{नव्या} \mr{पुलाचं} \mr{उद्घाटन} \mr{केलं}.
\end{quote}

\begin{quote}
\emph{mantryānnī navyā pulācaṁ udghāṭan kelaṁ.}
\end{quote}

\begin{quote}
The minister opened the new bridge.
\end{quote}

Masculine.

Respect takes the plural: \mr{मंत्री} \mr{आले}, the minister came.

\begin{cousinweb}
\textbf{\mr{मंत्री}} (\emph{mantrī}) is a tatsama, taken from Sanskrit \emph{mantrin}, \textquotedblleft{}a counsellor, a minister\textquotedblright{} (Monier-Williams). Marathi drops the final -n.
\end{cousinweb}

\begin{grammarlens}[title={What you've built}]
One more word for giving reasons and opinions.
\end{grammarlens}

\subsection*{Your turn}
\begin{itemize}
\raggedright
  \item \emph{Say it:} \emph{mantrī}
  \item \emph{Say it:} \emph{mantrī}, once more
  \item \emph{Say it:} \emph{mantryānnī navyā pulācaṁ udghāṭan kelaṁ.} — The minister opened the new bridge.
  \item \emph{Recall:} say the Marathi for administration, then the Marathi for parliament, then say \emph{mantrī} again
\end{itemize}

\begin{tcolorbox}[breakable,colback=teal!4,colframe=teal!35!black,title={Before you move on}]
What does \mr{मंत्री} mean? (\textquotedblleft{}A cabinet minister; a minister\textquotedblright{}.) Say it once more.
\end{tcolorbox}
\section[\texorpdfstring{\mr{नेता} (netā)}{नेता (netā)}]{\mr{नेता} (netā) — a leader}

\label{lesson:MR-C357-neta}



\begin{quote}
Before the new one: say the Marathi for parliament, then the Marathi for a cabinet minister.
\end{quote}

\subsection*{You'll want to know: \mr{नेता}}
\textbf{\mr{नेता}} — \emph{netā} — \textquotedblleft{}a leader\textquotedblright{}.

\begin{quote}
\mr{चांगला} \mr{नेता} \mr{लोकांचं} \mr{ऐकतो}.
\end{quote}

\begin{quote}
\emph{cāṅglā netā lokāñcaṁ aikto.}
\end{quote}

\begin{quote}
A good leader listens to people.
\end{quote}

Masculine. Plural: \mr{नेते}, \emph{nete}.

Plural \mr{नेते}; \mr{नेतृत्व}, leadership.

\begin{cousinweb}
\textbf{\mr{नेता}} (\emph{netā}) is a tatsama, taken from Sanskrit \emph{netṛ}, \textquotedblleft{}a leader, a guide\textquotedblright{} (Monier-Williams). The Sanskrit nominative is \emph{netā}.
\end{cousinweb}

\begin{grammarlens}[title={What you've built}]
One more word for giving reasons and opinions.
\end{grammarlens}

\subsection*{Your turn}
\begin{itemize}
\raggedright
  \item \emph{Say it:} \emph{netā}
  \item \emph{Say it:} \emph{netā}, once more
  \item \emph{Say it:} \emph{cāṅglā netā lokāñcaṁ aikto.} — A good leader listens to people.
  \item \emph{Recall:} say the Marathi for parliament, then the Marathi for a cabinet minister, then say \emph{netā} again
\end{itemize}

\begin{tcolorbox}[breakable,colback=teal!4,colframe=teal!35!black,title={Before you move on}]
What does \mr{नेता} mean? (\textquotedblleft{}A leader\textquotedblright{}.) Say it once more.
\end{tcolorbox}
\section[\texorpdfstring{\mr{नेमणे} (nemṇe)}{नेमणे (nemṇe)}]{\mr{नेमणे} (nemṇe) — to appoint}

\label{lesson:MR-C357-nemne}



\begin{quote}
Before the new one: say the Marathi for a cabinet minister, then the Marathi for a leader.
\end{quote}

\subsection*{You'll want to know: \mr{नेमणे}}
\textbf{\mr{नेमणे}} — \emph{nemṇe} — \textquotedblleft{}to appoint\textquotedblright{}.

\begin{quote}
\mr{कंपनीने} \mr{नवा} \mr{व्यवस्थापक} \mr{नेमला}.
\end{quote}

\begin{quote}
\emph{kampnīne navā vyavasthāpak nemlā.}
\end{quote}

\begin{quote}
The company appointed a new manager.
\end{quote}

Past, with -\mr{ने} (he …): \mr{त्याने} \mr{नेमलं}, \emph{tyāne nemlaṁ}.

X \mr{म्हणून} \mr{नेमणे}, to appoint as X; \mr{नेमणूक}, an appointment.

\begin{cousinweb}
\textbf{\mr{नेमणे}} (\emph{nemṇe}) is built on \mr{नेम}, \textquotedblleft{}a fixed rule, an aim\textquotedblright{}, usually traced to Sanskrit \emph{niyama}, \textquotedblleft{}restraint, fixed rule\textquotedblright{} (Monier-Williams).
\end{cousinweb}

\begin{grammarlens}[title={What you've built}]
One more word for giving reasons and opinions.
\end{grammarlens}

\subsection*{Your turn}
\begin{itemize}
\raggedright
  \item \emph{Say it:} \emph{nemṇe}
  \item \emph{Say it:} \emph{nemṇe}, once more
  \item \emph{Say it:} \emph{kampnīne navā vyavasthāpak nemlā.} — The company appointed a new manager.
  \item \emph{Recall:} say the Marathi for a cabinet minister, then the Marathi for a leader, then say \emph{nemṇe} again
\end{itemize}

\begin{tcolorbox}[breakable,colback=teal!4,colframe=teal!35!black,title={Before you move on}]
What does \mr{नेमणे} mean? (\textquotedblleft{}To appoint\textquotedblright{}.) Say it once more.
\end{tcolorbox}
\section[(dialogue)]{Looking back: from Accepting and Blaming to Who Governs}

\label{lesson:MR-R357-review-350-357}



\begin{quote}
Say the Marathi for to accept and to acknowledge.
\end{quote}

\subsection*{Your turn}
Say these aloud:

\begin{itemize}
\raggedright
  \item to accept, to acknowledge, to treat as valid, to blame, to accuse
  \item however, likewise, in fact, nevertheless, for example
  \item to mention, to state, to emphasize, especially, leading
  \item to confirm, to corroborate, to refute, actual, to react
  \item a citizen, the general public, common, population, rural
  \item voting, a representative, a candidate, canvassing, to be elected
  \item democracy, a majority, a minority, public opinion, to give a mandate
  \item administration, parliament, a cabinet minister, a leader, to appoint
\end{itemize}

\begin{tcolorbox}[breakable,colback=teal!4,colframe=teal!35!black,title={Before you move on}]
To accept? (\textbf{\mr{स्वीकारणे}}, \emph{svīkārṇe}.) To acknowledge? (\textbf{\mr{मान्य} \mr{करणे}}, \emph{mānya karṇe}.) To treat as valid? (\textbf{\mr{ग्राह्य} \mr{धरणे}}, \emph{grāhya dharṇe}.) To blame? (\textbf{\mr{दोष} \mr{देणे}}, \emph{doṣ deṇe}.) To accuse? (\textbf{\mr{आरोप} \mr{करणे}}, \emph{ārop karṇe}.) However? (\textbf{\mr{मात्र}}, \emph{mātra}.) Likewise? (\textbf{\mr{तसेच}}, \emph{tasec}.) In fact? (\textbf{\mr{किंबहुना}}, \emph{kimbahunā}.) Nevertheless? (\textbf{\mr{तथापि}}, \emph{tathāpi}.) For example? (\textbf{\mr{उदाहरणार्थ}}, \emph{udāharaṇārtha}.) To mention? (\textbf{\mr{उल्लेख} \mr{करणे}}, \emph{ullekh karṇe}.) To state? (\textbf{\mr{नमूद} \mr{करणे}}, \emph{namūd karṇe}.) To emphasize? (\textbf{\mr{भर} \mr{देणे}}, \emph{bhar deṇe}.) Especially? (\textbf{\mr{विशेषतः}}, \emph{viśeṣataḥ}.) Leading? (\textbf{\mr{प्रमुख}}, \emph{pramukh}.) To confirm? (\textbf{\mr{पुष्टी} \mr{करणे}}, \emph{puṣṭī karṇe}.) To corroborate? (\textbf{\mr{दुजोरा} \mr{देणे}}, \emph{dujorā deṇe}.) To refute? (\textbf{\mr{खंडन} \mr{करणे}}, \emph{khaṇḍan karṇe}.) Actual? (\textbf{\mr{वास्तविक}}, \emph{vāstavik}.) To react? (\textbf{\mr{प्रतिक्रिया} \mr{देणे}}, \emph{pratikriyā deṇe}.) A citizen? (\textbf{\mr{नागरिक}}, \emph{nāgarik}.) The general public? (\textbf{\mr{जनता}}, \emph{jantā}.) Common? (\textbf{\mr{सर्वसामान्य}}, \emph{sarvasāmānya}.) Population? (\textbf{\mr{लोकसंख्या}}, \emph{loksaṅkhyā}.) Rural? (\textbf{\mr{ग्रामीण}}, \emph{grāmīṇ}.) Voting? (\textbf{\mr{मतदान}}, \emph{matdān}.) A representative? (\textbf{\mr{प्रतिनिधी}}, \emph{pratinidhī}.) A candidate? (\textbf{\mr{उमेदवार}}, \emph{umedvār}.) Canvassing? (\textbf{\mr{प्रचार}}, \emph{pracār}.) To be elected? (\textbf{\mr{निवडून} \mr{येणे}}, \emph{nivḍūn yeṇe}.) Democracy? (\textbf{\mr{लोकशाही}}, \emph{lokśāhī}.) A majority? (\textbf{\mr{बहुमत}}, \emph{bahumat}.) A minority? (\textbf{\mr{अल्पसंख्याक}}, \emph{alpasaṅkhyāk}.) Public opinion? (\textbf{\mr{जनमत}}, \emph{janmat}.) To give a mandate? (\textbf{\mr{कौल} \mr{देणे}}, \emph{kaul deṇe}.) Administration? (\textbf{\mr{प्रशासन}}, \emph{praśāsan}.) Parliament? (\textbf{\mr{संसद}}, \emph{saṁsad}.) A cabinet minister? (\textbf{\mr{मंत्री}}, \emph{mantrī}.) A leader? (\textbf{\mr{नेता}}, \emph{netā}.) To appoint? (\textbf{\mr{नेमणे}}, \emph{nemṇe}.)
\end{tcolorbox}
